\documentclass[11pt]{amsart}

\usepackage[T1]{fontenc}
\usepackage{lmodern}
\usepackage{microtype}
\usepackage{mathtools,amssymb,amsthm}
\usepackage[hidelinks]{hyperref}

\hypersetup{
  pdftitle={Two non-bold wheel entries of Parida-Moura's Table 4 are equalities}
}

\newtheorem{theorem}{Theorem}
\newtheorem{lemma}[theorem]{Lemma}

\newcommand{\CFF}{\mathrm{CFF}}

\title[Two non-bold wheel entries are equalities]
{Two non-bold wheel entries of Parida--Moura's Table~4 are equalities:
 $t(W_{11})=t(W_{12})=8$}

\date{}

\subjclass[2020]{05D05, 05B40, 05C90}
\keywords{cover-free family, cover-free family on a graph, wheel, Sperner
family, exhaustive search}

\begin{document}

\begin{abstract}
For a graph $G$ on $n$ vertices let $t(G)$ be the least $t$ admitting a
$G$-cover-free family on the ground set $[t]$, in the sense of Parida and
Moura. Table~4 of their preprint tabulates $t(P_n)$, $t(C_n)$, $t(W_n)$ and
$t(K_n)$ for $n\le 12$ under the caption ``Upper bounds \dots\ where the bold
entries indicate exact values''; the two wheel entries $t(W_{11})\le 8$ and
$t(W_{12})\le 8$ are printed non-bold. We show that both are equalities,
\[
 t(W_{11})=t(W_{12})=8 .
\]
Table~4 leaves the exactness of these two non-bold upper bounds undetermined,
and no statement of that preprint asks whether they are exact; what we do is
supply the two matching lower bounds and so turn those two printed bounds into
equalities. The upper bounds come
from two families exhibited in full below. The lower bounds are exhaustive
non-existence statements: there is no $W_{11}$-$\CFF(7,11)$ and no
$W_{12}$-$\CFF(7,12)$, established by an exhaustive search over the whole
solution space.
\end{abstract}

\maketitle

\section{The definition and the two table entries}

Let $G=([1,n],E)$ be a graph. Following Parida and Moura
\cite[Definitions 3.1 and 4.1]{PM}, a family
$(B_1,\dots,B_n)$ of subsets of
$[1,t]$, with $B_i$ sitting on vertex $i$, is a \emph{$G$-cover-free family}
$G$-$\CFF(t,n)$ when
\begin{align}
 &B_a\setminus B_b\neq\varnothing
   \text{ and } B_b\setminus B_a\neq\varnothing
   &&(\{a,b\}\in E), \label{eq:S}\\
 &B_{i_0}\setminus(B_a\cup B_b)\neq\varnothing
   &&(\{a,b\}\in E,\ i_0\notin\{a,b\}),
   \label{eq:CF}
\end{align}
and
\[
 t(G)=\min\{\,t : \text{a } G\text{-}\CFF(t,|V(G)|) \text{ exists}\,\}.
\]
Throughout, $P_n$ and $C_n$ are the path and the cycle on $n$ vertices and
$W_n$ is the wheel on $n$ vertices: a hub joined to every vertex of a rim cycle
$C_{n-1}$. This is the convention of \cite{PM}, where $W_n$ is a universal
vertex joined to all vertices of a cycle of length $n-1$.

\medskip
The two entries at issue are in Table~4 of \cite{PM}, whose columns are $n$,
$t(P_n)$, $t(C_n)$, $t(W_n)$, $t(K_n)$. In the $t(W_n)$ column that table lists
a bold $7$ for $n=10$ and a plain $8$ for each of $n=11$ and $n=12$, under the
caption ``Upper bounds for $t(P_n)$, $t(C_n)$, $t(W_n)$, and $t(K_n)$, where
the bold entries indicate exact values.'' The table is introduced with the
words ``we display the known values for $t(P_n), t(C_n)$, and $t(W_n)$ or their
upper bounds for $n \leq 12$''. So by the source's own caption these two plain
entries stand there as upper bounds whose exactness is left undetermined; we
determine them below. No statement of \cite{PM} asks whether they are exact:
the source contains no conjecture, and its five Open Problem statements concern
other objects.

\begin{theorem}\label{thm:main}
$t(W_{11})=8$ and $t(W_{12})=8$.
\end{theorem}

Section~\ref{sec:upper} proves $\le$ by exhibiting the two families;
Section~\ref{sec:lower} proves $\ge$ by an exhaustive census.

\section{The upper bound: two exhibited families}\label{sec:upper}

The upper bound $8$ is already printed in \cite{PM}. The two paragraphs below
are elementary and are recorded only so that the two witnesses exhibited here
are self-contained; no novelty is claimed for them. We use the following
observations about \eqref{eq:S}--\eqref{eq:CF}.

\begin{lemma}\label{lem:anti}
Let $m\ge3$ and let $(B_1,\dots,B_m)$ be a $C_m$-$\CFF(t,m)$, with the cycle
$1,2,\dots,m,1$. Then $B_i\setminus B_j\neq\varnothing$ for all $i\neq j$; in
particular the $B_i$ are nonempty and pairwise distinct.
\end{lemma}

\begin{proof}
Fix $i\neq j$ and read indices modulo $m$. If $i\in\{j-1,j+1\}$ then $\{i,j\}$
is an edge and \eqref{eq:S} gives $B_i\setminus B_j\neq\varnothing$. Otherwise
$i\notin\{j,j+1\}$, so \eqref{eq:CF} applies to the edge $\{j,j+1\}$ with
$i_0=i$ and yields
$\varnothing\neq B_i\setminus(B_j\cup B_{j+1})\subseteq B_i\setminus B_j$.
\end{proof}

\begin{lemma}\label{lem:hub}
Let $m\ge3$, let $(B_1,\dots,B_m)$ be a $C_m$-$\CFF(t,m)$ on $[1,t]$, and set
$B_0=\{t+1\}$. Then $(B_0,B_1,\dots,B_m)$ is a $W_{m+1}$-$\CFF(t+1,m+1)$, where
the hub is vertex $0$ and the rim is the cycle $1,2,\dots,m,1$. Consequently
$t(W_{m+1})\le t(C_m)+1$.
\end{lemma}

\begin{proof}
Every $B_i$ with $i\ge1$ lies in $[1,t]$ and $B_0=\{t+1\}$ is disjoint from all
of them; all sets lie in $[1,t+1]$. There are two kinds of edge.

\emph{Hub edges $\{0,i\}$.} For \eqref{eq:S}:
$B_0\setminus B_i=\{t+1\}\neq\varnothing$, and
$B_i\setminus B_0=B_i\neq\varnothing$ by Lemma~\ref{lem:anti}. For
\eqref{eq:CF} with $i_0=k\notin\{0,i\}$:
$B_k\setminus(B_0\cup B_i)=B_k\setminus B_i\neq\varnothing$, again by
Lemma~\ref{lem:anti}.

\emph{Rim edges $\{i,i+1\}$.} Condition \eqref{eq:S} is inherited from the
$C_m$-$\CFF$. For \eqref{eq:CF} with $i_0=0$:
$B_0\setminus(B_i\cup B_{i+1})=\{t+1\}\neq\varnothing$; and with $i_0=k\ge1$,
$k\notin\{i,i+1\}$, it is inherited from the $C_m$-$\CFF$ as well.
\end{proof}

\begin{lemma}\label{lem:mono}
A $G$-$\CFF(t,n)$ is also a $G$-$\CFF(t+1,n)$. Hence if no $G$-$\CFF(t,n)$
exists then $t(G)\ge t+1$.
\end{lemma}

\begin{proof}
Neither \eqref{eq:S} nor \eqref{eq:CF} mentions the ground set; enlarging
$[1,t]$ to $[1,t+1]$ leaves every set, every difference and every union
unchanged, and $B_i\subseteq[1,t]\subseteq[1,t+1]$. The second sentence follows
because the set of $t$ admitting a $G$-$\CFF(t,n)$ is therefore an up-set.
\end{proof}

Here are the two objects. On the ground set $\{1,\dots,7\}$, indexed along the
rim cycle $v_1,v_2,\dots$, put
\begin{align}
 R_{10}&=\bigl(\{1\},\{2,3\},\{2,4\},\{2,5\},\{2,6\},
   \notag\\
       &\qquad\ \{6,7\},\{3,6\},\{3,5\},\{5,7\},\{4,7\}\bigr),
   \label{eq:R10}\\[2pt]
 R_{11}&=\bigl(\{1\},\{2,3,4\},\{2,3,5\},\{2,3,6\},\{2,7\},\{2,4,6\},
   \notag\\
       &\qquad\ \{4,5,6\},\{3,4,6\},\{3,4,7\},\{3,5,7\},\{5,6,7\}\bigr).
   \label{eq:R11}
\end{align}

\begin{lemma}\label{lem:witness}
$R_{10}$ is a $C_{10}$-$\CFF(7,10)$ and $R_{11}$ is a $C_{11}$-$\CFF(7,11)$.
\end{lemma}

\begin{proof}
Direct inspection of \eqref{eq:S}--\eqref{eq:CF}. For $R_{10}$ this is $10$
edges, the last of them the wrap edge $\{v_{10},v_1\}$, hence $20$ instances of
\eqref{eq:S} and $10\cdot 8=80$ instances of \eqref{eq:CF}; for $R_{11}$ it is
$11$ edges, $22$ instances of \eqref{eq:S} and $11\cdot 9=99$ of
\eqref{eq:CF}. Every one is a set difference between explicitly listed sets of
size at most $3$. The accompanying program performs all $221$ of them and
reports the number of violations, which is $0$ in both cases.
\end{proof}

\begin{proof}[Proof of $t(W_{11})\le 8$ and $t(W_{12})\le 8$]
Apply Lemma~\ref{lem:hub} with $t=7$ to $R_{10}$ and to $R_{11}$, which are
$C_{10}$- and $C_{11}$-cover-free families by Lemma~\ref{lem:witness}. With
$B_0=\{8\}$ this gives, on the ground set $\{1,\dots,8\}$, a
$W_{11}$-$\CFF(8,11)$
\[
 \bigl(\{8\}\,;\,R_{10}\bigr)
\]
and a $W_{12}$-$\CFF(8,12)$
\[
 \bigl(\{8\}\,;\,R_{11}\bigr),
\]
the hub carrying $\{8\}$ and the rim carrying \eqref{eq:R10} respectively
\eqref{eq:R11} in the printed order.
\end{proof}

\section{The lower bound: an exhaustive census}\label{sec:lower}

\begin{proof}[Proof of $t(W_{11})\ge 8$ and $t(W_{12})\ge 8$]
By Lemma~\ref{lem:mono} it suffices to show that no $W_{11}$-$\CFF(7,11)$ and
no $W_{12}$-$\CFF(7,12)$ exists. Each is a finite question: an assignment of a
subset of $[1,7]$ to each of $n$ vertices, so $(2^7)^n$ tuples, namely $2^{77}$
for $n=11$. It is decided as follows.

\emph{Symmetry.} Conditions \eqref{eq:S} and \eqref{eq:CF} are built from set
differences and unions only, so the whole constraint system is invariant under
the symmetric group $S_7$ acting on the ground set. The orbits of $S_7$ on
subsets of $[1,7]$ are exactly the size classes, so a solution exists if and
only if one exists in which the set on the first vertex of the search order is
one of the eight nested representatives
$\varnothing,\{1\},\{1,2\},\dots,\{1,\dots,7\}$. Restricting the first vertex
to those eight masks is therefore exact, not a heuristic.

\emph{Decomposition.} Fixing in addition the set on the second vertex of the
search order, arbitrarily among all $2^7=128$ subsets, splits that canonical
space into $8\times128=1024$ shards. They are pairwise disjoint, and their union
is the whole canonical space, because every canonical tuple has exactly one pair
of first two sets.

\emph{Enumeration.} Each shard is enumerated by depth-first search over the
remaining vertices, taking all $2^7$ subsets at each, and backtracking on
\eqref{eq:S} or \eqref{eq:CF} as soon as every vertex a condition mentions has
been assigned. No derived lemma is used as a prune, so an empty census is a
statement about \eqref{eq:S}--\eqref{eq:CF} alone.

\emph{Outcome.} All $1024$ of $1024$ expected shards return, and the totals are
\[
\begin{array}{lccc}
 & \text{solutions} & \text{shards} & \text{nodes}\\
 W_{11}\text{ at }t=7: & 0 & 1024/1024 & 14{,}002{,}112\\
 W_{12}\text{ at }t=7: & 0 & 1024/1024 & 14{,}015{,}792
\end{array}
\]
By Lemma~\ref{lem:mono}, $t(W_{11})\ge8$ and $t(W_{12})\ge8$.
\end{proof}

Theorem~\ref{thm:main} follows by combining the two sections.

\section{What was checked}\label{sec:status}

The accompanying program \texttt{verify.py}, whose recorded run is
\texttt{verify.output.txt}, checks the two printed witnesses
\eqref{eq:R10}--\eqref{eq:R11} and their hub extensions directly against
\eqref{eq:S}--\eqref{eq:CF} and re-runs the complete $t=7$ censuses for
$W_{11}$ and $W_{12}$, reproducing the shard, solution and node counts reported
above.

The upper bound is hand-checkable from \eqref{eq:R10}--\eqref{eq:R11} with $221$
elementary set differences; the lower bound rests on two exhaustive censuses,
which the program re-runs rather than replaces by a hand proof. Only the two
equalities of Theorem~\ref{thm:main} are established: the other non-bold cells
of Table~4 are neither proved nor used, no
$t(W_n)$ with $n\ge13$ is claimed, and no least-order claim is made --- nothing
here says $W_{11}$ is the smallest wheel whose table entry was left inexact.
The transcription and
interpretation of Table~4 are read by hand off the e-print of \cite{PM} rather
than produced by the computation.

\begin{thebibliography}{9}

\bibitem{PM}
P.~Parida and L.~Moura,
\emph{Cover-free families on graphs},
arXiv:2605.12634v1 [math.CO], 12 May 2026.
\href{https://arxiv.org/abs/2605.12634}{arXiv:2605.12634}.
Definition, table and numbering references are to v1.

\end{thebibliography}

\end{document}
